% GENERATED FILE. Edit canonical lessons, then run npm run generate:books.
% canonical-source-hash: fnv1a64:ee5ca9241a4b0601
% canonical-lessons: BN-C240-atkano, BN-C240-joma-deoya, BN-C240-jhap-deoya, BN-C240-doba, BN-C240-bhasa

\chapter{To block, To hand in, To dive, To sink, To float}
\label{ch:bn-to-block-to-hand-in-to-dive-to-sink-to-float}
\hlchaptermodality{240}

\begin{chapteropening}
\textbf{By the end of this chapter:} \emph{I can say to block, to hand in, to dive, to sink and to float.}

Complete the last lesson of chapter 240: I can say to block, to hand in, to dive, to sink and to float.
\end{chapteropening}

\section[\texorpdfstring{\bn{আটকানো} (āṭkāno)}{আটকানো (āṭkāno)}]{\bn{আটকানো} (āṭkāno) — to block, to hold up}

\label{lesson:BN-C240-atkano}



\begin{quote}
Before the new one: say the Bengali for to give back, then the Bengali for to remember.
\end{quote}

\subsection*{You'll want to know: \bn{আটকানো}}
\textbf{\bn{আটকানো}} — \emph{āṭkāno} — \textquotedblleft{}to block, to hold up\textquotedblright{}.

\begin{cousinweb}
\textbf{\bn{আটকানো}} (\emph{āṭkāno}) is a tadbhava, a word inherited and reshaped by sound change. It comes from \emph{*aṭṭakk}, \textquotedblleft{}obstruct, stop\textquotedblright{}, a form CDIAL reconstructs from its descendants (CDIAL 182). CDIAL lists Bengali \emph{āṭkāna}, \textquotedblleft{}to obstruct\textquotedblright{}.
\end{cousinweb}

\begin{grammarlens}[title={What you've built}]
One more word for this chapter.
\end{grammarlens}

\subsection*{Your turn}
\begin{itemize}
\raggedright
  \item \emph{Say it:} \emph{āṭkāno}
  \item \emph{Say it:} \emph{āṭkāno}, once more
  \item \emph{Say it:} say \emph{āṭkāno}
  \item \emph{Recall:} say the Bengali for to give back, then the Bengali for to remember, then say \emph{āṭkāno} again
\end{itemize}

\begin{tcolorbox}[breakable,colback=teal!4,colframe=teal!35!black,title={Before you move on}]
What does \bn{আটকানো} mean? (\textquotedblleft{}To block, to hold up\textquotedblright{}.) Say it once more.
\end{tcolorbox}
\section[\texorpdfstring{\bn{জমা} \bn{দেওয়া} (jômā deoyā)}{জমা দেওয়া (jômā deoyā)}]{\bn{জমা} \bn{দেওয়া} (jômā deoyā) — to hand in, to deposit}

\label{lesson:BN-C240-joma-deoya}



\begin{quote}
Before the new one: say the Bengali for to remember, then the Bengali for to block.
\end{quote}

\subsection*{You'll want to know: \bn{জমা} \bn{দেওয়া}}
\textbf{\bn{জমা} \bn{দেওয়া}} — \emph{jômā deoyā} — \textquotedblleft{}to hand in, to deposit\textquotedblright{}.

\begin{cousinweb}
\textbf{\bn{জমা} \bn{দেওয়া}} (\emph{jômā deoyā}) is \emph{jômā}, \textquotedblleft{}a deposit, a sum\textquotedblright{}, from Arabic \emph{jamʿ}, \textquotedblleft{}a gathering together\textquotedblright{}, which reached Bengali through Persian (Lane), + \emph{deoyā}, \textquotedblleft{}to give\textquotedblright{}, from Sanskrit \emph{dadāti}, \textquotedblleft{}gives\textquotedblright{} (CDIAL 6141).
\end{cousinweb}

\begin{grammarlens}[title={What you've built}]
One more word for this chapter.
\end{grammarlens}

\subsection*{Your turn}
\begin{itemize}
\raggedright
  \item \emph{Say it:} \emph{jômā deoyā}
  \item \emph{Say it:} \emph{jômā deoyā}, once more
  \item \emph{Say it:} say \emph{jômā deoyā}
  \item \emph{Recall:} say the Bengali for to remember, then the Bengali for to block, then say \emph{jômā deoyā} again
\end{itemize}

\begin{tcolorbox}[breakable,colback=teal!4,colframe=teal!35!black,title={Before you move on}]
What does \bn{জমা} \bn{দেওয়া} mean? (\textquotedblleft{}To hand in, to deposit\textquotedblright{}.) Say it once more.
\end{tcolorbox}
\section[\texorpdfstring{\bn{ঝাঁপ} \bn{দেওয়া} (jhā̃p deoyā)}{ঝাঁপ দেওয়া (jhā̃p deoyā)}]{\bn{ঝাঁপ} \bn{দেওয়া} (jhā̃p deoyā) — to dive}

\label{lesson:BN-C240-jhap-deoya}



\begin{quote}
Before the new one: say the Bengali for to block, then the Bengali for to hand in.
\end{quote}

\subsection*{You'll want to know: \bn{ঝাঁপ} \bn{দেওয়া}}
\textbf{\bn{ঝাঁপ} \bn{দেওয়া}} — \emph{jhā̃p deoyā} — \textquotedblleft{}to dive\textquotedblright{}.

\begin{cousinweb}
\textbf{\bn{ঝাঁপ} \bn{দেওয়া}} (\emph{jhāṁp deoyā}) is \emph{jhāṁp}, \textquotedblleft{}a sudden movement, a leap\textquotedblright{}, which CDIAL lists under \emph{*jhapp}, \textquotedblleft{}sudden movement\textquotedblright{}, a form it reconstructs (CDIAL 5336), + \emph{deoyā}, \textquotedblleft{}to give\textquotedblright{}, from Sanskrit \emph{dadāti}, \textquotedblleft{}gives\textquotedblright{} (CDIAL 6141).
\end{cousinweb}

\begin{grammarlens}[title={What you've built}]
One more word for this chapter.
\end{grammarlens}

\subsection*{Your turn}
\begin{itemize}
\raggedright
  \item \emph{Say it:} \emph{jhā̃p deoyā}
  \item \emph{Say it:} \emph{jhā̃p deoyā}, once more
  \item \emph{Say it:} say \emph{jhā̃p deoyā}
  \item \emph{Recall:} say the Bengali for to block, then the Bengali for to hand in, then say \emph{jhā̃p deoyā} again
\end{itemize}

\begin{tcolorbox}[breakable,colback=teal!4,colframe=teal!35!black,title={Before you move on}]
What does \bn{ঝাঁপ} \bn{দেওয়া} mean? (\textquotedblleft{}To dive\textquotedblright{}.) Say it once more.
\end{tcolorbox}
\section[\texorpdfstring{\bn{ডোবা} (ḍobā)}{ডোবা (ḍobā)}]{\bn{ডোবা} (ḍobā) — to sink, to drown}

\label{lesson:BN-C240-doba}



\begin{quote}
Before the new one: say the Bengali for to hand in, then the Bengali for to dive.
\end{quote}

\subsection*{You'll want to know: \bn{ডোবা}}
\textbf{\bn{ডোবা}} — \emph{ḍobā} — \textquotedblleft{}to sink, to drown\textquotedblright{}.

\begin{cousinweb}
\textbf{\bn{ডোবা}} (\emph{ḍobā}) is a tadbhava, a word inherited and reshaped by sound change. It comes from \emph{*ḍubb}, \textquotedblleft{}sink\textquotedblright{}, a form CDIAL reconstructs from its descendants (CDIAL 5561). CDIAL lists Bengali \emph{ḍubā}, \textquotedblleft{}to sink\textquotedblright{}, and takes the form as a reshuffling of Middle Indo-Aryan \emph{buḍḍaï}.
\end{cousinweb}

\begin{grammarlens}[title={What you've built}]
One more word for this chapter.
\end{grammarlens}

\subsection*{Your turn}
\begin{itemize}
\raggedright
  \item \emph{Say it:} \emph{ḍobā}
  \item \emph{Say it:} \emph{ḍobā}, once more
  \item \emph{Say it:} say \emph{ḍobā}
  \item \emph{Recall:} say the Bengali for to hand in, then the Bengali for to dive, then say \emph{ḍobā} again
\end{itemize}

\begin{tcolorbox}[breakable,colback=teal!4,colframe=teal!35!black,title={Before you move on}]
What does \bn{ডোবা} mean? (\textquotedblleft{}To sink, to drown\textquotedblright{}.) Say it once more.
\end{tcolorbox}
\section[\texorpdfstring{\bn{ভাসা} (bhāsā)}{ভাসা (bhāsā)}]{\bn{ভাসা} (bhāsā) — to float}

\label{lesson:BN-C240-bhasa}



\begin{quote}
Before the new one: say the Bengali for to dive, then the Bengali for to sink.
\end{quote}

\subsection*{You'll want to know: \bn{ভাসা}}
\textbf{\bn{ভাসা}} — \emph{bhāsā} — \textquotedblleft{}to float\textquotedblright{}.

\begin{cousinweb}
\textbf{\bn{ভাসা}} (\emph{bhāsā}) is a tadbhava, a word inherited from Sanskrit and reshaped by sound change. It comes from \emph{bhraśyati}, \textquotedblleft{}falls\textquotedblright{} (CDIAL 9654). CDIAL lists Bengali \emph{bhāsā}, \textquotedblleft{}to float, drift along\textquotedblright{}.
\end{cousinweb}

\begin{grammarlens}[title={What you've built}]
One more word for this chapter.
\end{grammarlens}

\subsection*{Your turn}
\begin{itemize}
\raggedright
  \item \emph{Say it:} \emph{bhāsā}
  \item \emph{Say it:} \emph{bhāsā}, once more
  \item \emph{Say it:} say \emph{bhāsā}
  \item \emph{Recall:} say the Bengali for to dive, then the Bengali for to sink, then say \emph{bhāsā} again
\end{itemize}

\begin{tcolorbox}[breakable,colback=teal!4,colframe=teal!35!black,title={Before you move on}]
What does \bn{ভাসা} mean? (\textquotedblleft{}To float\textquotedblright{}.) Say it once more.
\end{tcolorbox}
